\clearpage
\section{总结}
\subsection{可拓展之处}
预期实现风格权重在线下载：主机发送带模型版本、层结构、长度及CRC的参数包，先写入非活动存储区；RISC-V完成边界、位宽与完整性检查，等待引擎空闲后装载并原子切换，失败保留旧模型。当前仅实现从预置ROM加载，不宣称已有外部在线更新。另将继续补足十手势真人样本、改善弱光前景提取，并完成持续30fps实测及功耗、温升评估。

\subsection{心得体会}
本项目的主要难点是画质、存储、计算和交互共同约束下的系统一致性。减少参数并不必然降低全部资源；恢复全720p后，行缓存、访存和显示仍需处理完整像素。采用DW/PW和六通道降低运算量后，进一步的收益来自多上下文流水、列bank及总线重叠，不能只提高PLL频率。

训练方面，较强像素蒸馏会保留内容却削弱笔触，过强纹理又可能遮蔽物体。通过弱化像素回归、加入绘画Gram和局部纹理能量，当前梵高细笔触有所恢复，但有限感受野仍不能完成大范围语义重绘。评价应结合多内容图和真实摄像头，不能只看参考图或单张好例。

系统方面，完整EOF、最终DDR响应、输出槽元数据和安全换帧必须统一；任一环节提前发布都可能把算法问题与缓存撕裂混在一起。诊断计数使采集、计算和显示可分别观察，降低联调歧义。电气核对也必须以实际板卡原理图为依据，不能把DRC通过当成供电正确。

手势试验进一步说明了数据定义和真实泛化的重要性。用户的O形0、后三指3与拇指食指小指7并不等于数据集数字标签；自动标注也会把收拇指的手势认成7。四秒保持可以降低短时跳变，却无法消除持续误认。因此当前保留可解释的个人录入版本，未达标的通用模型不装进交付bit，后续需先补足正确数据和前景提取验证。

最后，交付不仅是bit文件，还需保留模型、整数参考、固件、工程配置和测量范围。本文将CAD与仿真结果、实际模型图及待完成实测分别说明，使后续更新能够沿着清楚的版本基线继续。
